(1) 单个位点大配分 \(\xi=1+e^{\beta(\mu+\epsilon)}\)，记 \(z=e^{\beta(\mu+\epsilon)}\)。独立位点使 \(\Xi=\xi^{N_s}\)，平均占据率 \(\boxed{\theta=z/(1+z)}\)；位点的平均吸附数 \(N_s\theta\)，满足 \(0<\theta<1\)。

(2) 占据变量是独立Bernoulli随机量，故 \(\langle N\rangle=N_s\theta\)、\(\operatorname{Var}N=N_s\theta(1-\theta)\)；Fano因子为 \(1-\theta\)。饱和 \(\theta\to1\) 时涨落趋零，和无限可占据的理想玻色粒子统计显著不同。

(3) 配置概率为单个位点 \(\{1-\theta,\theta\}\)，故 \(S_{\rm conf}=-N_sk_B[\theta\ln\theta+(1-\theta)\ln(1-\theta)]\)；平均吸附能为 \(U=-N_s\epsilon\theta\)。若 \(T\to0\) 且 \(\mu+\epsilon\ne0\)，占据趋0或1，熵趋零；恰在零温简并点时理想模型可出现残余配置简并。

(4) 代入气相化学势得 \(z=K(T)p\)，其中 \(K(T)=e^{\beta(\mu_0(T)+\epsilon)}/p_0\)，因此 \(\boxed{\theta=Kp/(1+Kp)}\)。如果位点之间有吸引、排斥或多个吸附态，则大配分不再因子化成独立Bernoulli因子，等温式和涨落均需重新计算。
